\documentclass[10pt,a4paper]{amsart}
\usepackage[margin=19mm]{geometry}
\usepackage{amsmath,amssymb,mathtools}
\usepackage[scr=rsfs,cal=euler]{mathalfa}
\usepackage{xfrac}
\usepackage[unicode,hidelinks,bookmarks=false]{hyperref}
\newcommand{\ie}{i.e.\ }
\newcommand{\cf}{cf.\ }
\newcommand{\resp}{respectively\ }
\newcommand{\Subsection}{\subsection}
\providecommand{\nobreakdash}{\nobreak\hbox{-}}
\setlength{\emergencystretch}{2em}
\multlinegap0pt
\usepackage{amssymb}
\usepackage{enumerate}
\usepackage{xcolor}
\usepackage{tikz}
\newtheorem{theorem}{Theorem}
\newcommand{\Chat}{\widehat{C}}
\newcommand{\Cov}{\mathbb{C}\mathrm{ov}}
\newcommand{\DK}{\Delta_k}
\newcommand{\EE}{\mathbb{E}}
\newcommand{\II}{\mathbb{I}}
\newcommand{\PP}{\mathbb{P}}
\newcommand{\SK}{\mathcal{S}_k}
\newcommand{\VV}{\mathbb{V}}
\newcommand{\Xhat}{\widehat{X}}
\title{Random discrete copulas}
\author{Damjana Kokol Bukov\v{s}ek, Bla\v{z} Moj\v{s}kerc, Nik Stopar}
\date{Source: arXiv:2603.13953v3 (CC BY 4.0)}
\begin{document}
\maketitle

\noindent\emph{Source and licence.} Random discrete copulas. Version arXiv:2603.13953v3 (CC BY 4.0), distributed by arXiv under the Creative Commons Attribution 4.0 International licence, http://creativecommons.org/licenses/by/4.0/, as recorded in the arXiv licence metadata for this exact version and shown on the abstract page https://arxiv.org/abs/2603.13953v3. Full licence notice: \texttt{licenses/arxiv-2603.13953-CC-BY-4.0.txt}.\par\smallskip

\noindent\textit{Editorial scope.} Counted unit: the article's theorem th:exk\_hat (source0.tex lines 577--582). The article's complete original proof of it is delivered with it (source0.tex lines 584--754, one \texttt{proof} environment of 171 lines). No mathematics is written, repaired or re-derived: the statement and the proof are the article's own lines, in the article's own notation, and no retained line is substituted or re-flowed.  Retained uncounted as the source-local context the proof needs: lines 267--272; lines 316--319; lines 356--359.  Nothing was omitted from any retained span, so the package carries no omission record.  No retained line contains a citation, so the wrapper carries no bibliography and no bibliography entry.  No retained line contains a figure environment, an image-inclusion command or drawing code, so no image is delivered and the package declares no asset dependency.  Editorial formatting only, in addition: an AMS \texttt{amsart} preamble that restates the article's own declarations and macros used by the retained lines, namely the environments \texttt{theorem} on separate counters, together with the standard packages those lines need.  The retained environments print in this document as Theorem 1 in order of appearance, whereas the article prints the same items with the numbering of its own full document.  The complete native source of the article as distributed in the arXiv e-print for this exact version is bundled unchanged as \texttt{sources/196363/source0.tex}.
\begin{align}
    |L_{ij\ell}|
    &= \binom{i}{\ell} \, \binom{j}{\ell} \, \ell! \, \binom{k-j}{i-\ell} \, (i-\ell)! \, \binom{k-i}{j-\ell} \, (j-\ell)! \,(k+\ell-i-j)! \nonumber\\
    &= \frac{i! \,j! \, \ell!\, (k-j)! \,(j-\ell)! \,(k-i)! \,(j-\ell)! \, (k+\ell-i-j)!}{\ell! \, (i-\ell)! \, \ell! \, (j-\ell)! \, (i-\ell)! \, (k+\ell-i-j)! \, (j-\ell)! \, (k+\ell-i-j)!} \label{eq:L_ijl_size}\\
    &= \frac{i! \, j! \, (k-i)! \, (k-j)!}{(i-\ell)! \, (j-\ell)! \, (k+\ell-i-j)! \, \ell!}. \nonumber
\end{align}

\begin{theorem} \label{th:exk}
    Let $k\ge 2$. For any $(u,v) \in \DK$ we have 
    $$ \EE[X_k(u,v)] = uv \qquad \text{ and } \qquad \VV[X_k(u,v)] = \frac{uv(1-u)(1-v)}{k-1}.$$
\end{theorem}

\begin{align}
    \EE\big[X_k(u,v)^2\big] &= \frac{i}{k^2 \, \binom{k}{j}} \binom{k-1}{j-1} + \frac{i(i-1)}{k^2 \binom{k}{j}} \binom{k-2}{j-2} \label{eq:E(X_k^2)}\\
    &= \frac{ij}{k^3} + \frac{i(i-1) \, j!\, (k-j)! \, (k-2)!}{k^2 \, k! \, (j-2)! \, (k-j)!} = \frac{ij}{k^3} + \frac{i(i-1)j(j-1)}{k^3 (k-1)}, \nonumber
\end{align}

\begin{theorem}\label{th:exk_hat}
    Let $k\ge 2$. For any $(u,v) \in \II^2$ let 
$$u = \tfrac{i}{k} + \tfrac{t}{k} \qquad \text{and} \qquad v = \tfrac{j}{k} + \tfrac{s}{k},$$
where $i=\lfloor uk \rfloor \in [k]_0$, $j=\lfloor vk \rfloor \in[k]_0$, $t = uk -i \in [0,1)$, and $s = vk -j \in [0,1)$. Then 
    $$ \EE[\Xhat_k(u,v)] = uv $$ and $$ \VV[\Xhat_k(u,v)] = \frac{1}{(k-1)}\Big(u(1-u)-\frac{t(1-t)}{k}\Big)\Big(v(1-v)-\frac{s(1-s)}{k}\Big).$$
\end{theorem}

\begin{proof}
Since $\Chat_\pi$ is a bilinear extension of $C_\pi$, for every $(u,v) \in \II^2$ the value $\Chat_\pi(u,v)$ is a convex combination of values $C_\pi(\tfrac{i}{k},\tfrac{j}{k})$, $C_\pi(\tfrac{i+1}{k},\tfrac{j}{k})$, $C_\pi(\tfrac{i}{k},\tfrac{j+1}{k})$, and $C_\pi(\tfrac{i+1}{k},\tfrac{j+1}{k})$
$$\Chat_\pi(u,v) = (1-t)(1-s)C_\pi(\tfrac{i}{k},\tfrac{j}{k}) + t(1-s)C_\pi(\tfrac{i+1}{k},\tfrac{j}{k}) + (1-t)sC_\pi(\tfrac{i}{k},\tfrac{j+1}{k}) +tsC_\pi(\tfrac{i+1}{k},\tfrac{j+1}{k}).$$
This implies that
\begin{align*}
    \Xhat_k(u,v) &= (1-t)(1-s)X_k(\tfrac{i}{k},\tfrac{j}{k}) + t(1-s)X_k(\tfrac{i+1}{k},\tfrac{j}{k}) \\
    & \qquad + (1-t)sX_k(\tfrac{i}{k},\tfrac{j+1}{k}) +tsX_k(\tfrac{i+1}{k},\tfrac{j+1}{k}).
\end{align*}
It follows from Theorem~\ref{th:exk} that the expected value of $\Xhat_k(u,v)$ is equal to
$$ \EE\big[\Xhat_k(u,v)\big] = (1-t)(1-s)\tfrac{ij}{k^2} + t(1-s)\tfrac{(i+1)j}{k^2} + (1-t)s\tfrac{i(j+1)}{k^2} +ts\tfrac{(i+1)(j+1)}{k^2}. $$
A straightforward algebraic manipulation reveals that the above expression simplifies to
$$ \EE\big[\Xhat_k(u,v)\big] =  \tfrac{(i+t)(j+s)}{k^2} = uv. $$

In order to compute the variance of $\Xhat_k(u,v)$, we need to find conditional distribution of  $X_k(\tfrac{i+1}{k},\tfrac{j}{k})$, $X_k(\tfrac{i}{k},\tfrac{j+1}{k})$, and $X_k(\tfrac{i+1}{k},\tfrac{j+1}{k})$ under the condition that $X_k(\tfrac{i}{k},\tfrac{j}{k}) = \tfrac{\ell}{k}$, since these random variables are not independent. 

Suppose that $X_k(\tfrac{i}{k},\tfrac{j}{k}) = \tfrac{\ell}{k}$. We may assume without loss of generality that $i<k$ and $j<k$. Recall the notation 
$$L_{ij\ell} = \{\pi \in \SK \mid C_\pi(\tfrac{i}{k},\tfrac{j}{k}) = \tfrac{\ell}{k}\}. $$
We split this set into five subsets
\begin{align*}
   L_{ij\ell}^1 &= \{\pi \in L_{ij\ell} \mid \pi(i+1)\le j, \pi^{-1}(j+1)\le i\}, \\
   L_{ij\ell}^2 &= \{\pi \in L_{ij\ell} \mid \pi(i+1)\le j, \pi^{-1}(j+1)\ge i+2\}, \\
   L_{ij\ell}^3 &= \{\pi \in L_{ij\ell} \mid \pi(i+1)\ge j+2, \pi^{-1}(j+1)\le i\}, \\
   L_{ij\ell}^4 &= \{\pi \in L_{ij\ell} \mid \pi(i+1)\ge j+2, \pi^{-1}(j+1)\ge i+2\}, \\
   L_{ij\ell}^5 &= \{\pi \in L_{ij\ell} \mid \pi(i+1) = j+1\}. 
\end{align*}
It is obvious that these sets are disjoint and their union is $L_{ij\ell}$. If $\pi \in L_{ij\ell}^1$, then 
$$|\{m \in [k] \mid m\le i, \pi(m) \le j\}| = \ell,$$
thus
\begin{align*}
    |\{m \in [k] \mid m\le i+1, \pi(m) \le j\}| &= \ell+1, \\
    |\{m \in [k] \mid m\le i, \pi(m) \le j+1\}| &= \ell+1,  \text{ and}\\
    |\{m \in [k] \mid m\le i+1, \pi(m) \le j+1\}| &= \ell+2. \\
\end{align*} 
It follows that
\begin{equation} \label{eq:Lijl1}
    \pi \in L_{ij\ell}^1 \Leftrightarrow C_\pi(\tfrac{i+1}{k},\tfrac{j}{k}) = \tfrac{\ell+1}{k},  C_\pi(\tfrac{i}{k},\tfrac{j+1}{k}) = \tfrac{\ell+1}{k},  C_\pi(\tfrac{i+1}{k},\tfrac{j+1}{k}) = \tfrac{\ell+2}{k}.
\end{equation} 
Similarly, we obtain
\begin{align*}
   \pi \in L_{ij\ell}^2 &\Leftrightarrow C_\pi(\tfrac{i+1}{k},\tfrac{j}{k}) = \tfrac{\ell+1}{k},  C_\pi(\tfrac{i}{k},\tfrac{j+1}{k}) = \tfrac{\ell}{k},  C_\pi(\tfrac{i+1}{k},\tfrac{j+1}{k}) = \tfrac{\ell+1}{k}, \\
   \pi \in L_{ij\ell}^3 &\Leftrightarrow C_\pi(\tfrac{i+1}{k},\tfrac{j}{k}) = \tfrac{\ell}{k},  C_\pi(\tfrac{i}{k},\tfrac{j+1}{k}) = \tfrac{\ell+1}{k},  C_\pi(\tfrac{i+1}{k},\tfrac{j+1}{k}) = \tfrac{\ell+1}{k}, \\
   \pi \in L_{ij\ell}^4 &\Leftrightarrow C_\pi(\tfrac{i+1}{k},\tfrac{j}{k}) = \tfrac{\ell}{k},  C_\pi(\tfrac{i}{k},\tfrac{j+1}{k}) = \tfrac{\ell}{k},  C_\pi(\tfrac{i+1}{k},\tfrac{j+1}{k}) = \tfrac{\ell}{k}, \\
   \pi \in L_{ij\ell}^5 &\Leftrightarrow C_\pi(\tfrac{i+1}{k},\tfrac{j}{k}) = \tfrac{\ell}{k},  C_\pi(\tfrac{i}{k},\tfrac{j+1}{k}) = \tfrac{\ell}{k},  C_\pi(\tfrac{i+1}{k},\tfrac{j+1}{k}) = \tfrac{\ell+1}{k}.
\end{align*}
We now compute the cardinality of $ L_{ij\ell}^1$.
To obtain a permutation $\pi\in L_{ij\ell}^1$, we first choose a subset $A \subseteq [i]$ containing $\ell$ elements and a subset $B \subseteq [j]$ containing $\ell$ elements and choose any bijective mapping from $A$ to $B$. Next we choose an element $a \in [i] \setminus A$ and map $a$ to $j+1$. Then we choose a subset $C \subseteq [k]\setminus [j+1]$ containing $i-\ell-1$ elements and choose any bijective mapping from $[i]\setminus (A\cup\{a\})$ to $C$. Next we choose an element $b \in [j] \setminus B$ and map $i+1$ to $b$. Then we choose a subset $D \subseteq [k]\setminus [i+1]$ containing $j-\ell-1$ elements and choose any bijective mapping from $D$ to $[j]\setminus (B\cup\{b\})$. Finally, we choose any bijective mapping from $[k]\setminus([i+1]\cup D)$ to $[k]\setminus([j+1]\cup C)$. Counting these choices, it follows that 
\begin{align*}
   |L_{ij\ell}^1| &= \binom{i}{\ell} \, \binom{j}{\ell} \, \ell! \, (i-\ell) \, \binom{k-j-1}{i-\ell-1} \, (i-\ell-1)! \, (j-\ell) \, \binom{k-i-1}{j-\ell-1} \, (j-\ell-1)! \,(k+\ell-i-j)! \\ 
   &= \frac{i! \,j! \, \ell!\, (i-\ell) \, (k-j-1)! \,(i-\ell-1)! \,(j-\ell) \, (k-i-1)! \,(j-\ell-1)! \, (k+\ell-i-j)!}{\ell! \, (i-\ell)! \, \ell! \, (j-\ell) \, (i-\ell-1)! \, (k+\ell-i-j)! \, (j-\ell-1)! \, (k+\ell-i-j)!} \\
   &= \frac{i! \, j! \, (k-i-1)! \, (k-j-1)!}{(i-\ell-1)! \, (j-\ell-1)! \, (k+\ell-i-j)! \, \ell!}. 
\end{align*}
We compute the cardinalities of the remaining sets in a similar way and obtain
\begin{align*}
   |L_{ij\ell}^2| &= \binom{i}{\ell} \, \binom{j}{\ell} \, \ell! \, \binom{k-j-1}{i-\ell} \, (i-\ell)! \, (j-\ell) \, \binom{k-i-1}{j-\ell-1} \, (j-\ell-1)! \,(k+\ell-i-j)! \\ 
   &= \frac{i! \, j! \, (k-i-1)! \, (k-j-1)!}{(i-\ell)! \, (j-\ell-1)! \, (k+\ell-i-j-1)! \, \ell!}, \\ 
   |L_{ij\ell}^3| &= \binom{i}{\ell} \, \binom{j}{\ell} \, \ell! \, (i-\ell) \, \binom{k-j-1}{i-\ell-1} \, (i-\ell-1)! \, \binom{k-i-1}{j-\ell} \, (j-\ell)! \,(k+\ell-i-j)! \\ 
   &= \frac{i! \, j! \, (k-i-1)! \, (k-j-1)!}{(i-\ell-1)! \, (j-\ell)! \, (k+\ell-i-j-1)! \, \ell!}, \\ 
   |L_{ij\ell}^4| &= \binom{i}{\ell} \, \binom{j}{\ell} \, \ell! \, \binom{k-j-1}{i-\ell} \, (i-\ell)! \, \binom{k-i-1}{j-\ell} \, (j-\ell)! \,(k+\ell-i-j-1)^2 \,(k+\ell-i-j-2)! \\ 
   &= \frac{i! \, j! \, (k-i-1)! \, (k-j-1)!}{(i-\ell)! \, (j-\ell)! \, (k+\ell-i-j-2)! \, \ell!}, \\ 
   |L_{ij\ell}^5| &= \binom{i}{\ell} \, \binom{j}{\ell} \, \ell! \, \binom{k-j-1}{i-\ell} \, (i-\ell)! \, \binom{k-i-1}{j-\ell} \, (j-\ell)! \,(k+\ell-i-j-1)! \\ 
   &= \frac{i! \, j! \, (k-i-1)! \, (k-j-1)!}{(i-\ell)! \, (j-\ell)! \, (k+\ell-i-j-1)! \, \ell!}.
\end{align*}
Using characterization \eqref{eq:Lijl1} and equation~\eqref{eq:L_ijl_size}, we compute the conditional probability
\begin{align*}
   p_{ij\ell}^1 &:=\PP\big[X_k(\tfrac{i+1}{k},\tfrac{j}{k}) = \tfrac{\ell+1}{k},  X_k(\tfrac{i}{k},\tfrac{j+1}{k}) = \tfrac{\ell+1}{k},  X_k(\tfrac{i+1}{k},\tfrac{j+1}{k}) = \tfrac{\ell+2}{k} \big| X_k(\tfrac{i}{k},\tfrac{j}{k}) = \tfrac{\ell}{k}\big] \\
   & \phantom{:}= \frac{|L_{ij\ell}^1|}{|L_{ij\ell}|} = \frac{\frac{i! \, j! \, (k-i-1)! \, (k-j-1)!}{(i-\ell-1)! \, (j-\ell-1)! \, (k+\ell-i-j)! \, \ell!}}{\frac{i! \, j! \, (k-i)! \, (k-j)!}{(i-\ell)! \, (j-\ell)! \, (k+\ell-i-j)! \, \ell!}} =\  \frac{(i-\ell) \, (j-\ell)}{(k-i) \, (k-j)}.
\end{align*}
Similarly,
\begin{align*}
   p_{ij\ell}^2 &:=\PP\big[X_k(\tfrac{i+1}{k},\tfrac{j}{k}) = \tfrac{\ell+1}{k},  X_k(\tfrac{i}{k},\tfrac{j+1}{k}) = \tfrac{\ell}{k}, X_k(\tfrac{i+1}{k},\tfrac{j+1}{k}) = \tfrac{\ell+1}{k} \big| X_k(\tfrac{i}{k},\tfrac{j}{k}) = \tfrac{\ell}{k}\big] \\
   & \phantom{:}= \frac{|L_{ij\ell}^2|}{|L_{ij\ell}|} = \frac{(j-\ell) \, (k+\ell-i-j)}{(k-i) \, (k-j)} , \\
   p_{ij\ell}^3 &:=\PP\big[X_k(\tfrac{i+1}{k},\tfrac{j}{k}) = \tfrac{\ell}{k},  X_k(\tfrac{i}{k},\tfrac{j+1}{k}) = \tfrac{\ell+1}{k},  X_k(\tfrac{i+1}{k},\tfrac{j+1}{k}) = \tfrac{\ell+1}{k} \big| X_k(\tfrac{i}{k},\tfrac{j}{k}) = \tfrac{\ell}{k}\big] \\
   & \phantom{:}= \frac{|L_{ij\ell}^3|}{|L_{ij\ell}|} = \frac{(i-\ell) \, (k+\ell-i-j)}{(k-i) \, (k-j)} , \\
   p_{ij\ell}^4 &:=\PP\big[X_k(\tfrac{i+1}{k},\tfrac{j}{k}) = \tfrac{\ell}{k},  X_k(\tfrac{i}{k},\tfrac{j+1}{k}) = \tfrac{\ell}{k},  X_k(\tfrac{i+1}{k},\tfrac{j+1}{k}) = \tfrac{\ell}{k} \big| X_k(\tfrac{i}{k},\tfrac{j}{k}) = \tfrac{\ell}{k}\big] \\
   & \phantom{:}= \frac{|L_{ij\ell}^4|}{|L_{ij\ell}|} = \frac{(k+\ell-i-j) \, (k+\ell-i-j-1)}{(k-i) \, (k-j)} , \\
   p_{ij\ell}^5 &:=\PP\big[X_k(\tfrac{i+1}{k},\tfrac{j}{k}) = \tfrac{\ell}{k},  X_k(\tfrac{i}{k},\tfrac{j+1}{k}) = \tfrac{\ell}{k},  X_k(\tfrac{i+1}{k},\tfrac{j+1}{k}) = \tfrac{\ell+1}{k} \big| X_k(\tfrac{i}{k},\tfrac{j}{k}) = \tfrac{\ell}{k}\big] \\
   & \phantom{:}= \frac{|L_{ij\ell}^5|}{|L_{ij\ell}|} = \frac{(k+\ell-i-j)}{(k-i) \, (k-j)} .
\end{align*}
If $\pi \in L_{ij\ell}^1$, condition \eqref{eq:Lijl1} implies that $\Chat_\pi(u,v) =  \frac{\ell+t+s}{k}$. If $\pi \in L_{ij\ell}^2$, we have $\Chat_\pi(u,v) =  \frac{\ell+t}{k}$, if $\pi \in L_{ij\ell}^3$, we obtain $\Chat_\pi(u,v) =  \frac{\ell+s}{k}$, if $\pi \in L_{ij\ell}^4$, we have $\Chat_\pi(u,v) =  \frac{\ell}{k}$, and, finally, if $\pi \in L_{ij\ell}^5$, we get $\Chat_\pi(u,v) =  \frac{\ell+ts}{k}$.
So, the conditional distribution of random variable $\Xhat_k(u,v)$ under the condition that $X_k(\tfrac{i}{k},\tfrac{j}{k}) = \tfrac{\ell}{k}$ is
$$\big(\Xhat_k(u,v) \big| X_k(\tfrac{i}{k},\tfrac{j}{k}) = \tfrac{\ell}{k}\big) \sim  \begin{pmatrix}
    \frac{\ell+t+s}{k} & \frac{\ell+t}{k} & \frac{\ell+s}{k} & \frac{\ell}{k} & \frac{\ell+ts}{k} \\[2mm]
    p_{ij\ell}^1 & p_{ij\ell}^2 & p_{ij\ell}^3 & p_{ij\ell}^4 & p_{ij\ell}^5
\end{pmatrix}.$$
We remark that in this formula some of the values may collapse into a single one. For example, when $s=t$, we have $\frac{\ell+t}{k}=\frac{\ell+s}{k}$, and this value is attained with probability $p_{ij\ell}^2+p_{ij\ell}^3$.
We denote the sets
$$M_{st\ell} = \{\tfrac{\ell+t+s}{k}, \tfrac{\ell+t}{k}, \tfrac{\ell+s}{k}, \tfrac{\ell}{k}, \tfrac{\ell+ts}{k}\} \qquad \text{and} \qquad M_{uv} = \bigcup_{\ell=\max \{0,  i+j-k \}}^{\min \{i,j \}} M_{st\ell}.$$
We have proved that $\Xhat_k(u,v)$ is a discrete random variable with range in $M_{uv}$. Note that in the case $t > 0$, $s > 0$, $t\ne s$, and $t+s \ne 1$ each set $M_{st\ell}$ contains five different elements and the sets $M_{st\ell}$ are pairwise disjoint, so that 
\begin{align*}
    \PP\big[\Xhat_k(u,v) = \tfrac{\ell+t+s}{k}\big] &= p_{ij\ell}^1\, \PP[X_k(\tfrac{i}{k},\tfrac{j}{k}) = \tfrac{\ell}{k}], \\
    \PP\big[\Xhat_k(u,v) = \tfrac{\ell+t}{k}\big] &= p_{ij\ell}^2\, \PP[X_k(\tfrac{i}{k},\tfrac{j}{k}) = \tfrac{\ell}{k}], \\
    \PP\big[\Xhat_k(u,v) = \tfrac{\ell+s}{k}\big] &= p_{ij\ell}^3\, \PP[X_k(\tfrac{i}{k},\tfrac{j}{k}) = \tfrac{\ell}{k}], \\
    \PP\big[\Xhat_k(u,v) = \tfrac{\ell}{k}\big] &= p_{ij\ell}^4\, \PP[X_k(\tfrac{i}{k},\tfrac{j}{k}) = \tfrac{\ell}{k}], \text{ and}\\
    \PP\big[\Xhat_k(u,v) = \tfrac{\ell+ts}{k}\big] &= p_{ij\ell}^5\, \PP[X_k(\tfrac{i}{k},\tfrac{j}{k}) = \tfrac{\ell}{k}]. 
\end{align*}
If $0 < t = s \ne \tfrac12$ each set $M_{st\ell}$ contains four different elements and the sets $M_{st\ell}$ are still pairwise disjoint, so that in this case
$$\PP\big[\Xhat_k(u,v) = \tfrac{\ell+t}{k}\big] = (p_{ij\ell}^2 + p_{ij\ell}^3)\, \PP[X_k(\tfrac{i}{k},\tfrac{j}{k}) = \tfrac{\ell}{k}],$$
and the other probabilities are the same as above. If $t > 0, s > 0$, $t\ne s$, and $t+s=1$ the sets are not disjoint, namely $M_{st\ell} \cap M_{st(\ell+1)} = \{\tfrac{\ell+1}{k}\}$, so that in this case
$$\PP\big[\Xhat_k(u,v) = \tfrac{\ell+1}{k}\big] = p_{ij\ell}^1\, \PP[X_k(\tfrac{i}{k},\tfrac{j}{k}) = \tfrac{\ell}{k}] + p_{ij(\ell+1)}^4\, \PP[X_k(\tfrac{i}{k},\tfrac{j}{k}) = \tfrac{\ell+1}{k}],$$
and the other probabilities are the same as above. If $t=0$ or $s=0$ or $t=s=\tfrac12$ even more possible values of $\Xhat_k(u,v)$ collapse into a single value. 

Furthermore, conditional probabilities of the values of $X_k(\tfrac{i+1}{k},\tfrac{j}{k})$, $X_k(\tfrac{i}{k},\tfrac{j+1}{k})$, and $X_k(\tfrac{i+1}{k},\tfrac{j+1}{k})$ under the condition that $X_k(\tfrac{i}{k},\tfrac{j}{k})  = \tfrac{\ell}{k}$ are
\begin{align*}
   \PP\big[X_k(\tfrac{i+1}{k},\tfrac{j}{k}) = \tfrac{\ell}{k} \big| X_k(\tfrac{i}{k},\tfrac{j}{k}) = \tfrac{\ell}{k}\big] 
   &= p_{ij\ell}^3 + p_{ij\ell}^4 + p_{ij\ell}^5 = \frac{k+\ell-i-j}{k-i}, \\
   \PP\big[X_k(\tfrac{i+1}{k},\tfrac{j}{k}) = \tfrac{\ell+1}{k} \big| X_k(\tfrac{i}{k},\tfrac{j}{k}) = \tfrac{\ell}{k}\big] 
   &= p_{ij\ell}^1 + p_{ij\ell}^2 = \frac{j-\ell}{k-i}, \\
   \PP\big[X_k(\tfrac{i}{k},\tfrac{j+1}{k}) = \tfrac{\ell}{k} \big| X_k(\tfrac{i}{k},\tfrac{j}{k}) = \tfrac{\ell}{k}\big] 
   &= p_{ij\ell}^2 + p_{ij\ell}^4 + p_{ij\ell}^5 = \frac{k+\ell-i-j}{k-j}, \\
   \PP\big[X_k(\tfrac{i}{k},\tfrac{j+1}{k}) = \tfrac{\ell+1}{k} \big| X_k(\tfrac{i}{k},\tfrac{j}{k}) = \tfrac{\ell}{k}\big] 
   &= p_{ij\ell}^1 + p_{ij\ell}^3 = \frac{i-\ell}{k-j}, \\
   \PP\big[X_k(\tfrac{i+1}{k},\tfrac{j+1}{k}) = \tfrac{\ell}{k} \big| X_k(\tfrac{i}{k},\tfrac{j}{k}) = \tfrac{\ell}{k}\big] 
   &= p_{ij\ell}^4 = \frac{(k+\ell-i-j) \, (k+\ell-i-j-1)}{(k-i) \, (k-j)}, \\
   \PP\big[X_k(\tfrac{i+1}{k},\tfrac{j+1}{k}) = \tfrac{\ell+1}{k} \big| X_k(\tfrac{i}{k},\tfrac{j}{k}) = \tfrac{\ell}{k}\big] 
   &= p_{ij\ell}^2 + p_{ij\ell}^3 + p_{ij\ell}^5 = \frac{(k+\ell-i-j) \, (i+j+1-2\ell)}{(k-i) \, (k-j)}, \\
   \PP\big[X_k(\tfrac{i+1}{k},\tfrac{j+1}{k}) = \tfrac{\ell+2}{k} \big| X_k(\tfrac{i}{k},\tfrac{j}{k}) = \tfrac{\ell}{k}\big] 
   &= p_{ij\ell}^1 = \frac{(i-\ell) \, (j-\ell)}{(k-i) \, (k-j)}.
\end{align*}

We are now able to compute the covariances between random variables $X_k(\tfrac{i}{k},\tfrac{j}{k})$, $X_k(\tfrac{i+1}{k},\tfrac{j}{k})$, $X_k(\tfrac{i}{k},\tfrac{j+1}{k})$, and $X_k(\tfrac{i+1}{k},\tfrac{j+1}{k})$.
We have
\begin{align*}
    \EE\big[X_k(\tfrac{i}{k},\tfrac{j}{k})X_k(\tfrac{i+1}{k},\tfrac{j}{k})\big] &= \sum_{\ell=\max \{0,  i+j-k \}}^{\min \{i,j \}} \left(\frac{\ell^2}{k^2}\cdot \PP\big[X_k(\tfrac{i}{k},\tfrac{j}{k}) = \tfrac{\ell}{k}, X_k(\tfrac{i+1}{k},\tfrac{j}{k}) = \tfrac{\ell}{k}\big] \right.\\
    & \left.\qquad + \frac{\ell(\ell+1)}{k^2}\cdot \PP\big[X_k(\tfrac{i}{k},\tfrac{j}{k}) = \tfrac{\ell}{k}, X_k(\tfrac{i+1}{k},\tfrac{j}{k})= \tfrac{\ell+1}{k}\big]\right)\\
    &= \sum_{\ell=\max \{0,  i+j-k \}}^{\min \{i,j \}} \frac{\ell^2}{k^2}\cdot \PP\big[X_k(\tfrac{i}{k},\tfrac{j}{k}) = \tfrac{\ell}{k}\big] \\
    & \qquad + \sum_{\ell=\max \{0,  i+j-k \}}^{\min \{i,j \}} \frac{\ell}{k^2}\cdot \PP\big[X_k(\tfrac{i}{k},\tfrac{j}{k}) = \tfrac{\ell}{k}, X_k(\tfrac{i+1}{k},\tfrac{j}{k})= \tfrac{\ell+1}{k}\big].
\end{align*}
The first sum is equal to $\EE\big[X_k(\tfrac{i}{k},\tfrac{j}{k})^2\big]$, while in the second sum we use the fact that the probability is equal to $(p_{ij\ell}^1 + p_{ij\ell}^2)\PP\big[X_k(\tfrac{i}{k},\tfrac{j}{k}) = \tfrac{\ell}{k}\big]$.
Thus, we obtain
\begin{align*}
    \EE\big[X_k(\tfrac{i}{k},\tfrac{j}{k})X_k(\tfrac{i+1}{k},\tfrac{j}{k})\big] &= \EE\big[X_k(\tfrac{i}{k},\tfrac{j}{k})^2\big] + \sum_{\ell=\max \{0,  i+j-k \}}^{\min \{i,j \}} \frac{\ell}{k^2}\cdot \frac{j-\ell}{k-i}\cdot \PP\big[X_k(\tfrac{i}{k},\tfrac{j}{k}) = \tfrac{\ell}{k}\big]\\
    &= \EE\big[X_k(\tfrac{i}{k},\tfrac{j}{k})^2\big] + \frac{j}{k(k-i)}\sum_{\ell=\max \{0,  i+j-k \}}^{\min \{i,j \}} \frac{\ell}{k}\cdot \PP\big[X_k(\tfrac{i}{k},\tfrac{j}{k}) = \tfrac{\ell}{k}\big] \\
    & \qquad - \frac{1}{k-i}\sum_{\ell=\max \{0,  i+j-k \}}^{\min \{i,j \}} \frac{\ell^2}{k^2}\cdot \PP\big[X_k(\tfrac{i}{k},\tfrac{j}{k}) = \tfrac{\ell}{k}\big].
\end{align*}
Again, the first sum is equal to $\EE\big[X_k(\tfrac{i}{k},\tfrac{j}{k})\big]$, and the second to $\EE\big[X_k(\tfrac{i}{k},\tfrac{j}{k})^2\big]$, so by~\eqref{eq:E(X_k^2)} we have
\begin{align*}
    \EE\big[X_k(\tfrac{i}{k},\tfrac{j}{k})X_k(\tfrac{i+1}{k},\tfrac{j}{k})\big] &= \EE\big[X_k(\tfrac{i}{k},\tfrac{j}{k})^2\big] + \frac{j}{k(k-i)}\EE\big[X_k(\tfrac{i}{k},\tfrac{j}{k})\big]
    - \frac{1}{k-i}\EE\big[X_k(\tfrac{i}{k},\tfrac{j}{k})^2\big] \\
    &=  \frac{k-i-1}{k-i} \cdot \left(\frac{ij}{k^3} + \frac{i(i-1)j(j-1)}{k^3 (k-1)}\right) + \frac{j}{k(k-i)}\cdot \frac{ij}{k^2} \\
    &=  \frac{ij(k-i-1+ij)}{k^3(k-1)}.
\end{align*}
This implies that
\begin{align*}
    \Cov[X_k(\tfrac{i}{k},\tfrac{j}{k}), X_k(\tfrac{i+1}{k},\tfrac{j}{k})] &= \EE\big[X_k(\tfrac{i}{k},\tfrac{j}{k})X_k(\tfrac{i+1}{k},\tfrac{j}{k})\big] - \EE[X_k(\tfrac{i}{k},\tfrac{j}{k})]\EE[X_k(\tfrac{i+1}{k},\tfrac{j}{k})] \\
    &=  \frac{ij(k-i-1+ij)}{k^3(k-1)} - \frac{ij}{k^2}\cdot\frac{(i+1)j}{k^2} \\
    &=  \frac{ij(k-i-1)(k-j)}{k^4(k-1)}. 
\end{align*}
Similarly, we obtain
$$\Cov[X_k(\tfrac{i}{k},\tfrac{j}{k}), X_k(\tfrac{i}{k},\tfrac{j+1}{k})] = \frac{ij(k-i)(k-j-1)}{k^4(k-1)}, $$
$$\Cov[X_k(\tfrac{i}{k},\tfrac{j}{k}), X_k(\tfrac{i+1}{k},\tfrac{j+1}{k})] = \frac{ij(k-i-1)(k-j-1)}{k^4(k-1)}, $$
and
$$\Cov[X_k(\tfrac{i+1}{k},\tfrac{j}{k}), X_k(\tfrac{i}{k},\tfrac{j+1}{k})] = \frac{ij(k-i-1)(k-j-1)}{k^4(k-1)}. $$
Finally, expanding the variance of $\Xhat_k(u,v)$, we obtain
\begin{align*}
    \VV\big[\Xhat_k(u,v)\big] &= (1-t)^2(1-s)^2\cdot\VV[X_k(\tfrac{i}{k},\tfrac{j}{k})] + t^2(1-s)^2\cdot\VV[X_k(\tfrac{i+1}{k},\tfrac{j}{k})] \\
    & \qquad + (1-t)^2s^2\cdot\VV[X_k(\tfrac{i}{k},\tfrac{j+1}{k})]+t^2s^2 \cdot\VV[X_k(\tfrac{i+1}{k},\tfrac{j+1}{k})] \\
    & \qquad + 2(1-t)t(1-s)^2\cdot\Cov[X_k(\tfrac{i}{k},\tfrac{j}{k}), X_k(\tfrac{i+1}{k},\tfrac{j}{k})]  \\
    & \qquad + 2(1-t)ts^2\cdot\Cov[X_k(\tfrac{i}{k},\tfrac{j+1}{k}), X_k(\tfrac{i+1}{k},\tfrac{j+1}{k})] \\
    & \qquad + 2(1-t)^2(1-s)s\cdot\Cov[X_k(\tfrac{i}{k},\tfrac{j}{k}), X_k(\tfrac{i}{k},\tfrac{j+1}{k})]  \\
    & \qquad + 2t^2(1-s)s\cdot\Cov[X_k(\tfrac{i+1}{k},\tfrac{j}{k}), X_k(\tfrac{i+1}{k},\tfrac{j+1}{k})] \\
    & \qquad + 2(1-t)t(1-s)s\cdot\Cov[X_k(\tfrac{i}{k},\tfrac{j}{k}), X_k(\tfrac{i+1}{k},\tfrac{j+1}{k})]  \\
    & \qquad + 2(1-t)t(1-s)s\cdot\Cov[X_k(\tfrac{i+1}{k},\tfrac{j}{k}), X_k(\tfrac{i}{k},\tfrac{j+1}{k})].
\end{align*}
Plugging in the expressions for the (co-) variances and simplifying, we get
\begin{align*}
    \VV\big[\Xhat_k(u,v)\big]
    &=\frac{\big((i+t)(k-i-t)-kt(1-t)\big)\big((j+s)(k-j-s)-ks(1-s)\big)}{k^4(k-1)} \\
    &= \frac{1}{(k-1)}\Big(u(1-u)-\frac{t(1-t)}{k}\Big)\Big(v(1-v)-\frac{s(1-s)}{k}\Big),
\end{align*}
which finishes the proof.
\end{proof}

\end{document}
